\scanpage{75}\mspage{74}

Cela \unsure{résultera|résulte} du Corollaire qui suit, et de $C\unsure{5|\deleted{6}5}$

\uline{Corollaire} (Propriété de la relation $\leq$ !) Soient $F$ une figure, $K$ une sous-figure \deleted{qui} telle que $K\neq F$. Alors il existe une \deleted{\unsure{str}} figure $G$ (une multistrate, si on veut) qui est \deleted{\unsure{donc|des}} disjointe de $K$, mais non de $F$. En d'autres termes, $\operatorname{supp} K\neq\operatorname{supp} F$.

Ce n'est pas ``formel'' en termes des axiomes $C1 : C4$ sur $(\mathcal F,\leq)$, Supposons p.ex. que $\mathcal F$ admette un plus grand élément, donc il est défini comme l'ens de \uline{toutes} les parties fermées de l'ens ordonné $M$, et ``disjoint'' signifie ici $A\cap B=\phi$ sans plus (la relation de compatibilité $R$ est la relation tjrs vraie). Mais si $F$ est une partie fermée, $K\subset F$ une sous partie fermée \inserted{$\neq F$}, il n'est pas clair qu'on puisse trouver une partie fermée $G$ telle que $G\cap\deleted{F}K=\phi$, $G\cap F\neq\phi$. \deleted{P.ex. $M\unsure{=|-}[0,N]$, $\unsure{F|R}$ avec \unsure{rel\textsuperscript{n}} \unsure{d'ordre|l'ordre} \unsure{(?)|ord} habituel, $F=M$, $K=[0,1]$,} \inserted{Si p.ex. $M$ a un plus petit \unsure{élément|élém\textsuperscript{t}|élt} $0$,} Alors pour \uline{toute} partie fermée non vide, les $G$ tels que $F$ et $G$ disjoints sont les $G$ vides, donc ça ne change pas si on remplace $F$ par $\deleted{\unsure{F|R}}K\in F$, \unsure{tant|\deleted{tant}} que $K$ n'est pas vid\unsure{e}
